% !TeX root = ../main.tex
\section{Conclusion}
\label{sec:conclusion}

Guitar tablature transcription has moved from frame-wise CNNs on acoustic
guitar to modular systems with learned fingering and technique stages, and to
generative end-to-end models. In parallel, large sequence-to-sequence
transcribers trained on real multi-instrument data now handle full mixes and
distorted electric guitar, but stop at notes. For the use case of
transcribing a guitar part from a commercial song, the practical state of the
art is a modular pipeline: a band-split or hybrid separator, a guitar note
model trained on many tones (or a multi-instrument transcriber), a neural beat
tracker, and a learned or search-based fingering stage. The most promising
direction is to close the split between the two lines of work, by giving
mix-level transcribers string and fret outputs, and to build a benchmark of
real recordings with reference tablatures so that progress on the actual task
can be measured.
